\documentclass[11pt]{article}
\usepackage[margin=1in]{geometry}
\usepackage{enumitem}
% =============================================================================
% Preambles/header.tex — shared LaTeX/Beamer preamble for this project
%
% Use from a Beamer lecture by prepending:
%     \input{header}
% and compiling with TEXINPUTS=../Preambles:$TEXINPUTS (the /compile-latex
% skill does this automatically).
%
% PALETTE CONTRACT. Color names defined here MUST match the names in
% Quarto/theme-template.scss so Beamer and Quarto renderings use the same
% palette. The scripts/check-palette-sync.sh script enforces this.
%
% To customize: edit the HEX values below AND the matching $variable in
% Quarto/theme-template.scss, then run scripts/check-palette-sync.sh.
% =============================================================================

% -----------------------------------------------------------------------------
% Engine detection + encoding
%
% Our /compile-latex skill uses XeLaTeX, where inputenc is unnecessary and
% can produce encoding warnings. Only load inputenc under pdfTeX.
% -----------------------------------------------------------------------------
\usepackage{iftex}
\ifPDFTeX
  \usepackage[utf8]{inputenc}
\fi

\usepackage{amsmath, amssymb, amsthm}
\usepackage{booktabs}
\usepackage{graphicx}

% -----------------------------------------------------------------------------
% PALETTE — mirrors Quarto/theme-template.scss variable names.
% Load xcolor and define colors BEFORE anything that references them
% (notably hyperref, hypersetup, and the Beamer color assignments).
% -----------------------------------------------------------------------------
\usepackage{xcolor}

\definecolor{primary-blue}{HTML}{012169}      % headings, accents
\definecolor{primary-gold}{HTML}{B9975B}      % emphasis, borders
\definecolor{highlight-yellow}{HTML}{F2A900}  % markers, alerts
\definecolor{light-bg}{HTML}{E8EDF5}          % title / section backgrounds
\definecolor{jet}{HTML}{1A1A1A}               % body text

% Semantic colors (used in diagrams and annotations)
\definecolor{positive}{HTML}{15803D}          % good / observed / identified
\definecolor{negative}{HTML}{B91C1C}          % bad / problematic / bias
\definecolor{neutral}{HTML}{525252}           % reference / context

% Highlight accents (match .hi-* classes in the SCSS)
\definecolor{hi-slate}{HTML}{314F4F}
\definecolor{hi-green}{HTML}{2E8B57}
\definecolor{hi-red}{HTML}{C41E3A}

% Semantic aliases so skill templates and snippets can write
% \textcolor{accent}{...} without hard-coding "primary-blue".
\colorlet{accent}{primary-blue}
\colorlet{accent2}{primary-gold}

% -----------------------------------------------------------------------------
% Hyperref — loaded AFTER xcolor + palette so link colors resolve.
% -----------------------------------------------------------------------------
\usepackage{hyperref}
\hypersetup{colorlinks=true, linkcolor=primary-blue, urlcolor=primary-blue,
            citecolor=primary-blue, pdfborder={0 0 0}}

% -----------------------------------------------------------------------------
% Course code — set once per deck (after \input{header}, before \begin{document})
% via \coursecode{CS401}. Shown in the footer on every slide so a deck's course
% is identifiable without opening the title page. Empty by default.
% -----------------------------------------------------------------------------
\makeatletter
\newcommand{\coursecode}[1]{\gdef\@coursecodetext{#1}}
\gdef\@coursecodetext{}
\makeatother

% -----------------------------------------------------------------------------
% Beamer theme assignments (only apply under Beamer).
%
% We detect Beamer via \@ifclassloaded{beamer}, which is the documented,
% non-internal way. This must be wrapped in \makeatletter/\makeatother
% because \@ifclassloaded contains an @-catcode character.
% -----------------------------------------------------------------------------
\makeatletter
\@ifclassloaded{beamer}{%
  \usefonttheme{professionalfonts}
  \setbeamercolor{structure}{fg=primary-blue}
  \setbeamercolor{frametitle}{fg=primary-blue}
  \setbeamercolor{title}{fg=primary-blue}
  \setbeamercolor{subtitle}{fg=primary-gold}
  \setbeamercolor{author}{fg=primary-blue}
  \setbeamercolor{institute}{fg=neutral}
  \setbeamercolor{date}{fg=primary-gold}
  \setbeamercolor{itemize item}{fg=primary-blue}
  \setbeamercolor{itemize subitem}{fg=primary-gold}
  \setbeamercolor{alerted text}{fg=primary-gold}
  \setbeamercolor{block title}{fg=white, bg=primary-blue}
  \setbeamercolor{block body}{bg=light-bg}
  % Semantic block variants: block=definition/concept (blue), exampleblock=
  % worked example (green/positive), alertblock=key takeaway or caution (gold).
  % Keep to <=2 colored boxes per slide across ALL three variants (INV-7).
  \setbeamercolor{block title example}{fg=white, bg=positive}
  \setbeamercolor{block body example}{bg=light-bg}
  \setbeamercolor{block title alerted}{fg=white, bg=primary-gold}
  \setbeamercolor{block body alerted}{bg=light-bg}

  % Minimal footer: course code (left) + page X / N (right), no navigation clutter
  \setbeamertemplate{navigation symbols}{}
  \setbeamertemplate{footline}{%
    \color{neutral}\footnotesize\hspace{0.7em}\@coursecodetext\hfill
    \insertframenumber/\inserttotalframenumber\hspace{0.7em}\vspace{0.3em}}
}{}
\makeatother

% -----------------------------------------------------------------------------
% TikZ setup — libraries the snippet gallery uses
% -----------------------------------------------------------------------------
\usepackage{tikz}
\usetikzlibrary{arrows.meta, positioning, calc, decorations.pathreplacing,
                fit, shapes.geometric, backgrounds}

% Shared TikZ styles — snippets and hand-written diagrams can pick these up
% via `\begin{tikzpicture}[every node/.style={...}]` or by naming specific
% styles (e.g., `\node[dag-node] (X) at (0,0) {X};`).
\tikzset{
  % Node styles for causal DAGs and similar
  dag-node/.style = {
    draw=primary-blue, fill=light-bg, rounded corners=2pt,
    minimum width=1.2cm, minimum height=0.9cm, align=center, font=\small
  },
  decision-node/.style = {
    draw=primary-gold, fill=white, diamond, aspect=1.8,
    minimum width=1.8cm, minimum height=1.2cm, align=center, font=\small,
    inner sep=1pt
  },
  % Edge styles — observed vs counterfactual
  observed-edge/.style = {-{Latex[length=2.5mm]}, thick, draw=jet},
  counterfactual-edge/.style = {-{Latex[length=2.5mm]}, thick, draw=neutral, dashed},
  confound-edge/.style = {-{Latex[length=2.5mm]}, thick, draw=negative, dashed},
  % Data-point styles for plots
  observed-dot/.style = {fill=primary-blue, draw=primary-blue, circle, inner sep=1.2pt},
  counterfactual-dot/.style = {fill=white, draw=neutral, circle, inner sep=1.2pt}
}

% -----------------------------------------------------------------------------
% Convenience macros
% -----------------------------------------------------------------------------
\newcommand{\muted}[1]{\textcolor{neutral}{#1}}
\newcommand{\key}[1]{\textcolor{primary-gold}{\textbf{#1}}}
\newcommand{\good}[1]{\textcolor{positive}{#1}}
\newcommand{\bad}[1]{\textcolor{negative}{#1}}

% Standout frame macro: full-bleed dark background for section transitions.
% Beamer-only (uses \begin{frame}/\setbeamercolor/\usebeamerfont) -- do not
% call from an article-class file (e.g. Notes/<CODE>/*-notes.tex). \muted,
% \key, \good, \bad, and \coursecode above are plain xcolor/gdef and are
% safe to use from either class; verified when Notes/article-class support
% was added.
% Expand: \transitionslide{Your transition title}
\newcommand{\transitionslide}[1]{%
  \begingroup
  \setbeamercolor{background canvas}{bg=primary-blue}
  \begin{frame}[plain]
    \centering
    \vfill
    {\usebeamerfont{title}\color{white}\Large #1\par}
    \vfill
  \end{frame}
  \endgroup
}

% Section divider frame (Beamer-only), an alternative to \transitionslide with a
% darker two-line style: near-black (jet) background, a small white label line
% above a larger gold title, and a thin gold rule along the bottom edge.
% Expand: \sectiondivider{Section 2}{Pointers}
\newcommand{\sectiondivider}[2]{%
  \begingroup
  \setbeamercolor{background canvas}{bg=jet}
  \begin{frame}[plain]
    \vspace*{3.0cm}
    {\color{white}\ttfamily\large #1\par}
    \vspace{0.5cm}
    {\color{primary-gold}\LARGE #2\par}
    \begin{tikzpicture}[remember picture, overlay]
      \draw[primary-gold, line width=2.5pt]
        ($(current page.south west)+(0,0.1)$) -- ($(current page.south east)+(0,0.1)$);
    \end{tikzpicture}
  \end{frame}
  \endgroup
}

% -----------------------------------------------------------------------------
% Channel watermark — "Concepts That Click" (YouTube).
%
% Article-class files (CompetitiveExam Q/A sets, Notes, Assignments) get a
% light diagonal draft watermark on every page plus a centered footer naming
% the channel. Beamer decks get a subtle TikZ background watermark instead
% (the footline already carries the course code). Centralized here so every
% MCQ PDF is branded without per-file edits.
% -----------------------------------------------------------------------------
\makeatletter
\@ifclassloaded{article}{%
  \usepackage{draftwatermark}
  \SetWatermarkText{Concepts That Click}
  \SetWatermarkScale{1}
  \SetWatermarkFontSize{2cm}
  \SetWatermarkLightness{0.86}
  \SetWatermarkAngle{45}
  \usepackage{fancyhdr}
  \pagestyle{fancy}
  \fancyhf{}
  \fancyfoot[C]{\footnotesize\textcolor{neutral}{Concepts That Click~|~YouTube: \href{https://www.youtube.com/@concepts.thatclick}{@concepts.thatclick}}}
  \renewcommand{\headrulewidth}{0pt}
  \renewcommand{\footrulewidth}{0pt}
  \fancypagestyle{plain}{%
    \fancyhf{}%
    \fancyfoot[C]{\footnotesize\textcolor{neutral}{Concepts That Click~|~YouTube: \href{https://www.youtube.com/@concepts.thatclick}{@concepts.thatclick}}}%
    \renewcommand{\headrulewidth}{0pt}%
    \renewcommand{\footrulewidth}{0pt}%
  }
}{}
\@ifclassloaded{beamer}{%
  \setbeamertemplate{background}{%
    \begin{tikzpicture}[remember picture, overlay]
      \node[rotate=45, opacity=0.055, align=center]
        at (current page.center)
        {\Huge\textcolor{neutral}{Concepts That Click}};
    \end{tikzpicture}%
  }
}{}
\makeatother 
\coursecode{JKSSB}
\renewcommand{\thesection}{49.\arabic{section}}
\newtheorem{example}{Example}[section]
\newenvironment{definitionbox}[1]{%
  \par\noindent\textbf{Definition (#1).}\ \itshape
}{\par\vspace{0.3em}}
\title{Computer Generations \& History --- Detailed Lecture Notes}
\author{JKSSB Computer Awareness}
\date{\today}

\begin{document}
\maketitle

\section{What a computer generation represents}
Computer generations are conventional historical groupings based on the
dominant hardware technology or design development associated with a period.
They help organize the major changes in computer hardware, programming, and
use. They are not exact universal year boundaries: a date range can differ
between accounts, so the technology association is the safer cue for an
objective question.

\begin{definitionbox}{Computer generation}
A conventional historical grouping of computers associated with a dominant
technology or design development.
\end{definitionbox}

\begin{center}
\begin{tabular}{p{0.22\textwidth}p{0.68\textwidth}}
\toprule
\textbf{Term} & \textbf{Meaning in this lesson}\\
\midrule
Vacuum tube & An early electronic component used in first-generation computers; these systems were large and produced substantial heat.\\
Transistor & A smaller electronic component that replaced vacuum tubes in the conventional second-generation account.\\
Integrated circuit (IC) & A compact circuit incorporating multiple electronic components; the key third-generation association.\\
Microprocessor & A processor implemented as an integrated chip; the key fourth-generation association.\\
VLSI & Very Large Scale Integration: integrating a very large number of circuit components on a chip.\\
Machine language & Instructions expressed in the low-level codes directly used by the machine.\\
Assembly language & A symbolic, low-level notation for machine operations; it must be translated before execution.\\
High-level language & A programming language using more abstract statements than machine code or assembly.\\
Operating system & System software that manages the computer's operation and provides services to programs and users.\\
Multiprogramming & Keeping multiple programs available so the system can organize processor use while one program is waiting.\\
\bottomrule
\end{tabular}
\end{center}

The language and system-software terms help explain what changed beyond the
hardware labels. ``Machine language'' identifies the first generation's
programming cue. Assembly and early high-level languages appear with the
second-generation association. Operating systems and multiprogramming
develop with the third-generation systems described here. These terms are
not alternate names for the hardware generation; they are developments
associated with it.

\section{The first four technology associations}
\begin{center}
\begin{tabular}{p{0.20\textwidth}p{0.27\textwidth}p{0.43\textwidth}}
\textbf{Generation} & \textbf{Main association} & \textbf{Exam cue}\\
\hline
First & Vacuum tubes & Large, heat-producing systems; machine-language programming.\\
Second & Transistors & Smaller and more reliable; assembly and early high-level languages.\\
Third & Integrated circuits (ICs) & Multiple electronic components integrated
in compact circuits; smaller, more reliable systems; operating systems and
multiprogramming develop.\\
Fourth & Microprocessors & VLSI integration; personal computers become
widespread.\\
\end{tabular}
\end{center}

The progression is a change in how electronic components are built and
combined. Vacuum tubes made early systems large and heat-producing.
Transistors replaced those tubes and supported smaller, more reliable
systems. Integrated circuits placed multiple components together in a
compact circuit, continuing the trend toward reduced size and improved
reliability. In the fourth-generation association, a microprocessor brings
processor functions together on a chip; VLSI integration and the spread of
personal computers are the accompanying cues.

The sequence is the important memory structure:
\[
\text{vacuum tubes}\ \longrightarrow\ \text{transistors}\
\longrightarrow\ \text{integrated circuits}\ \longrightarrow\
\text{microprocessors}.
\]
The most common near-neighbour error is to switch the third and fourth
generations. An integrated circuit is the third-generation cue; the
microprocessor is the fourth-generation cue.

\begin{example}
An item asks which generation is associated with integrated circuits. Trace
the sequence: vacuum tubes are first, transistors second, ICs third, and
microprocessors fourth. Therefore the answer is the third generation. If
the stem instead names a microprocessor, the conventional answer is fourth.
\end{example}

\section{The fifth-generation association}
The conventional fifth-generation association commonly includes Artificial
Intelligence (AI), parallel processing, and natural-language systems. This
is a broad category, not a claim that every modern computer has one
universally agreed start date or a single defining design. AI applications
can run on systems that are not otherwise classified by one precise
``fifth-generation'' boundary. For an exam stem that explicitly uses this
generation scheme, match the conventional association and avoid treating
the label as a sharply dated technical standard.

\begin{example}
If the stem asks which broad generation is commonly linked to AI and
parallel processing in the conventional classification, select the fifth.
If it asks for the microprocessor generation, select the fourth instead;
``modern'' alone is not a sufficient clue to erase the distinction.
\end{example}

\section{India's PARAM 8000 milestone}
For the named national milestone, C-DAC's official high-performance
computing history identifies:
\href{https://cdac.in/index.aspx?id=hpc_ss_supercomputing_systems}
{C-DAC, HPC Supercomputing Systems}.

\begin{definitionbox}{C-DAC}
Centre for Development of Advanced Computing.
\end{definitionbox}

\begin{definitionbox}{PARAM 8000}
An Indian supercomputer developed by C-DAC and unveiled in 1991; C-DAC
identifies it as India's first supercomputer.
\end{definitionbox}

Keep the three parts of this fact separate. C-DAC is the organization;
PARAM 8000 is the system; 1991 is the unveiling year. This milestone
belongs to India's high-performance-computing history. It is not a
generation label, and it should not be used to redefine the conventional
technology sequence.

\begin{example}
For a matching item with ``Centre for Development of Advanced Computing,''
select C-DAC. For a prompt asking for India's first supercomputer, select
PARAM 8000. If the prompt asks when the system was unveiled, use 1991.
These answers refer to related but different parts of the same fact.
\end{example}

\section{Exam retrieval and common confusions}
Use two retrieval directions. Given the technology, recall the generation;
given the generation, recall the technology. Then isolate named milestones
that sit outside the generation sequence.
\begin{itemize}
  \item Vacuum tubes $\leftrightarrow$ first generation.
  \item Transistors $\leftrightarrow$ second generation.
  \item Integrated circuits $\leftrightarrow$ third generation.
  \item Microprocessors $\leftrightarrow$ fourth generation.
  \item AI and parallel processing $\leftrightarrow$ common fifth-generation
    description.
  \item C-DAC / PARAM 8000 / 1991 $\leftrightarrow$ separate Indian
    supercomputing milestone.
\end{itemize}

Three distinctions prevent most errors in this topic:
\begin{enumerate}
  \item Do not swap the adjacent third- and fourth-generation cues: ICs are
    third; microprocessors are fourth.
  \item Do not treat ``fifth generation'' as a single precisely dated
    hardware boundary. AI, parallel processing, and natural-language
    systems are the conventional association used in this lesson.
  \item Do not mix a generation label with the PARAM 8000 fact. The latter
    identifies a named Indian supercomputer, its developer, and its
    unveiling year.
\end{enumerate}

\subsection*{Exercises}
\begin{enumerate}[label=\arabic*.]
  \item Match vacuum tubes, transistors, integrated circuits, and
    microprocessors to the first four generations.
  \item Explain why the presence of an IC inside a modern microprocessor
    does not change the conventional microprocessor/fourth-generation
    association.
  \item Expand C-DAC and identify PARAM 8000's significance.
  \item Why should fifth-generation date boundaries be treated cautiously?
  \item Distinguish a generation milestone from the PARAM 8000 milestone.
\end{enumerate}

\subsection*{Solutions}
\begingroup\small
\begin{enumerate}[label=\arabic*.]
  \item First: vacuum tubes; second: transistors; third: integrated
    circuits; fourth: microprocessors.
  \item The generation scheme uses the dominant historical milestone as its
    label. A later microprocessor can contain ICs without changing the
    conventional fourth-generation association.
  \item C-DAC is Centre for Development of Advanced Computing. PARAM 8000
    was developed by C-DAC and is identified as India's first
    supercomputer.
  \item Generation boundaries are broad historical groupings and do not
    have one universally fixed transition date; technology associations
    are the more stable recall cue.
  \item A generation milestone classifies a historical technology period.
    PARAM 8000 is a named Indian supercomputer and national HPC milestone,
    not one of the generation labels.
\end{enumerate}
\endgroup
\end{document}
